\documentclass[12pt]{article}
\usepackage[utf8]{inputenc}
\usepackage[T1]{fontenc}
\usepackage{hyperref}
\title{Western Xwla Gbe Language Resources for K-12 Curriculum Development}
\author{Compiled by Claude Sonnet 5.5 (claude-sonnet-5-5) for Andr\'as Kornai}
\date{}
\begin{document}
\maketitle

\section{Introduction}
Western Xwla Gbe (ISO 639-3 xwl; also called Phla or Xwla) is a Gbe language of the Niger-Congo family, spoken in Benin and Togo. It belongs to the Phla--Pher\'a cluster. Speaker estimates differ: about 94,000 in one SIL-based dataset and 160,000 (2018) in Wikipedia. Wikipedia treats Eastern (gbx) and Western (xwl) as separate codes. The searches found very little educational material.

\section{Primary Educational Resources}
No primers, textbooks or school programmes for Western Xwla Gbe were found in the searches.

\section{Linguistic and Reference Materials}
\begin{itemize}
\item Wikipedia, Phla language: \url{https://en.wikipedia.org/wiki/Phla_language}
\item Wikipedia, Phla--Pher\'a languages: \url{https://en.wikipedia.org/wiki/Phla%E2%80%93Pher%C3%A1_languages}
\item Numeral system data (MPI-EVA Chan numerals database): \url{https://lingweb.eva.mpg.de/channumerals/Gbe-Western-Xwla.htm}
\end{itemize}

\section{Religious and Community Resources}
\begin{itemize}
\item Scripture Earth listing for xwl: \url{https://www.scriptureearth.org/00eng.php?iso=xwl}
\item Joshua Project profile, which reports scripture portions and few Christian resources in the language: \url{https://joshuaproject.net/people_groups/15150/BN}
\end{itemize}

\section{Digital and Audio Resources}
No audio, video or software resources were verified in the searches.

\section{Curriculum Development Context}
The searches found no evidence of use in schools, and little written material. A curriculum would need new development work, probably starting from the wider Gbe literacy tradition and community orthography decisions, which these searches did not cover.
\end{document}